\documentclass[11pt,a4paper]{article}

\usepackage[utf8]{inputenc}
\usepackage[T1]{fontenc}
\usepackage{amsmath,amssymb,amsthm}
\usepackage{geometry}
\usepackage{hyperref}
\usepackage{booktabs}

\geometry{margin=1in}

\theoremstyle{plain}
\newtheorem{theorem}{Theorem}
\newtheorem{lemma}[theorem]{Lemma}
\newtheorem{corollary}[theorem]{Corollary}
\newtheorem{proposition}[theorem]{Proposition}
\newtheorem{conjecture}[theorem]{Conjecture}

\theoremstyle{definition}
\newtheorem{observation}[theorem]{Observation}
\newtheorem{remark}[theorem]{Remark}
\newtheorem{definition}[theorem]{Definition}

\title{Block Reduction and Mutual Exclusion\\
in the Collatz Corona\\
\large Version 2}
\author{Oriol Corcoll Arias\\
\small Working draft}
\date{May 2026}

\begin{document}

\maketitle

\begin{abstract}
For the accelerated Collatz map, the cyclic component is governed by the
integrality condition $(2^x - 3^y) \mid C(y, \vec\sigma)$, where
$C(y, \vec\sigma) = \sum_{k=0}^{y-1} 3^{y-1-k} 2^{\sigma_k}$ is the cycle
constant of an admissible exponent pattern $\vec\sigma$. The set
$\mathcal{C}(x, y) = \{C(y, \vec\sigma)\}$ is the \emph{Collatz corona} of
Belaga--Mignotte. Across $1{,}769{,}706{,}522$ patterns enumerated over 22
parameter pairs, no non-trivial integrality is observed.

The main contribution of this version is a rigorous algebraic mechanism that
explains a substantial fraction of this avoidance. We prove a
\emph{block reduction lemma}: when $\gcd(x, y) = g \geq 2$ and $\vec\sigma$
is the $r$-fold repetition (for $r \mid g$, $r \geq 2$) of a base block $B$
of length $b = y/r$ and weight $a = x/r$, the cycle constant satisfies the
exact integer identity
\[
C(y, \vec\sigma) = C(b, B) \cdot \frac{2^x - 3^y}{2^a - 3^b}.
\]
Combined with the algebraic factorisation
$2^x - 3^y = (2^a - 3^b) \, \Phi_{a, b, r}$,
this yields a recursive reduction of the cyclic obstruction within periodic
patterns: cycle integrality at $(x, y)$ for $\vec\sigma = B^r$ is equivalent
to cycle integrality at $(a, b)$ for $B$, modulo a coprimality condition
that holds in all empirically tested cases.

The lemma forces algebraic divisibility of $C(y, \vec\sigma)$ by every prime
factor of $\Phi_{a, b, r}$, regardless of the choice of $B$. This explains
the empirical \emph{mutual exclusion} between divisibility events at distinct
prime divisors of $2^x - 3^y$ recorded in the first version of this note: the
factor $\Phi_{a, b, r}$ divides $C$ automatically, while the complementary
factor $2^a - 3^b$ divides $C$ only via the strictly smaller cycle equation
on $B$.

Patterns admit a natural three-fold partition: \emph{periodic} (covered by
the lemma), \emph{sporadic} (primitive period $y$ within reducible parameter
pairs), and \emph{irreducible} (parameter pairs with $\gcd(x, y) = 1$). The
lemma applies only to periodic patterns. The conjecture is shown to reduce
to its statement on sporadic patterns of reducible pairs and on irreducible
parameter pairs.

The note does not prove the absence of non-trivial Collatz cycles. It
isolates an algebraic mechanism, identifies the location of the unsolved
part of the obstruction, and reformulates the cyclic component of Collatz
as a problem about irreducible base cases.
\end{abstract}


\section{Setup and prior work}

The accelerated Collatz map on odd positive integers is
\[
T(n) = \frac{3n + 1}{2^{v_2(3n + 1)}}.
\]
A periodic orbit with $y$ odd steps and $x$ total halvings is encoded by an
admissible exponent pattern
$\vec\sigma = (\sigma_0, \ldots, \sigma_{y-1})$,
$0 = \sigma_0 < \sigma_1 < \cdots < \sigma_{y-1} < x$.
The classical cycle equation \cite{bohm1978,lagarias1985} is
\begin{equation}\label{eq:cycle}
n_0 (2^x - 3^y) = C(y, \vec\sigma) := \sum_{k=0}^{y-1} 3^{y-1-k} 2^{\sigma_k}.
\end{equation}
A non-trivial cycle requires $2^x > 3^y$ and the integrality condition
$(2^x - 3^y) \mid C(y, \vec\sigma)$.

Equivalently, write $\vec\sigma$ as a positive composition $a_1, \ldots, a_y \ge 1$
with $\sum a_k = x$, where $a_k = \sigma_k - \sigma_{k-1}$ (and $a_y = x - \sigma_{y-1}$). Let
\[
\Sigma(x, y) = \{\text{admissible patterns}\}, \qquad |\Sigma(x, y)| = \binom{x-1}{y-1}.
\]
Following \cite{belaga1998}, the integer $2^x - 3^y$ is the \emph{Collatz number}
associated with $(x, y)$, and the finite set
$\mathcal{C}(x, y) = \{C(y, \vec\sigma) : \vec\sigma \in \Sigma(x, y)\}$
is the \emph{Collatz corona}.

The trivial cycle $1 \to 1$ corresponds to $a_1 = \cdots = a_y = 2$, requiring
$x = 2y$.

Existing rigorous results on cycles approach the problem via Diophantine
estimates and linear forms in logarithms; see Steiner~\cite{steiner1977},
Eliahou~\cite{eliahou1993}, Simons--de Weger~\cite{sdw2005}, Simons~\cite{simons2008},
Hercher~\cite{hercher2022}, with the current bound $m \ge 92$ on the number
of local minima of any non-trivial cycle. Generalised settings appear in
Belaga~\cite{belaga2003}, Belaga--Mignotte~\cite{belaga2006},
Holden~\cite{holden2011}. Recent work of Dhiman--Pandey~\cite{dhiman2026}
shows that the unrestricted cycle divisibility predicate
$\mathcal{D}_y = \{(x, C) \in \mathbb{N}^2 : (2^x - 3^y) \mid C\}$
is not Presburger-definable.


\section{The block reduction lemma}\label{sec:block}

The main rigorous result of this version isolates the algebraic structure of
the Collatz corona on patterns admitting a periodic block decomposition.

\begin{lemma}[Block reduction]\label{lem:block}
Let $r \ge 1$, and write $x = r a$, $y = r b$ with $a, b \ge 1$. Let
$B = (b_1, \ldots, b_b)$ be a positive composition of $a$, and let
$\vec\sigma = (B, B, \ldots, B)$ denote the $r$-fold concatenation. Then
\begin{equation}\label{eq:block-identity}
C(y, \vec\sigma) = C(b, B) \cdot \sum_{j=0}^{r-1} 2^{j a} \, 3^{(r-1-j) b}.
\end{equation}
\end{lemma}

\begin{proof}
The partial sums of $\vec\sigma$ are
$\sigma_{j b + i} = j a + S_i$,
where $S_0 = 0, S_1, \ldots, S_{b-1}$ are the partial sums of $B$, for
$0 \le j \le r - 1$ and $0 \le i \le b - 1$. Substituting into \eqref{eq:cycle}:
\begin{align*}
C(y, \vec\sigma)
&= \sum_{j=0}^{r-1} \sum_{i=0}^{b-1} 3^{r b - 1 - j b - i} \, 2^{j a + S_i} \\
&= \sum_{j=0}^{r-1} 2^{j a} \, 3^{(r - 1 - j) b}
   \sum_{i=0}^{b-1} 3^{b - 1 - i} \, 2^{S_i}
 = C(b, B) \cdot \sum_{j=0}^{r-1} 2^{j a} \, 3^{(r - 1 - j) b}. \qedhere
\end{align*}
\end{proof}

\begin{proposition}[Algebraic factorisation]\label{prop:factorisation}
For every $r, a, b \ge 1$ with $2^a \neq 3^b$,
\begin{equation}\label{eq:diff-factorisation}
2^x - 3^y = (2^a - 3^b) \, \Phi_{a, b, r},
\qquad
\Phi_{a, b, r} := \sum_{j=0}^{r-1} 2^{j a} \, 3^{(r - 1 - j) b},
\end{equation}
where $x = r a$, $y = r b$.
\end{proposition}

\begin{proof}
Apply the standard identity $A^r - B^r = (A - B) \sum_{j=0}^{r-1} A^{r-1-j} B^j$
to $A = 2^a$, $B = 3^b$. Reindexing, $\sum_{j=0}^{r-1} (2^a)^{r-1-j} (3^b)^j
= \sum_{j=0}^{r-1} 2^{(r-1-j) a} 3^{j b}$. The result \eqref{eq:diff-factorisation}
follows after relabelling the summation index.
\end{proof}

\begin{corollary}[Forced divisibility]\label{cor:forced}
Let $r, a, b \ge 1$ with $2^a > 3^b$, and let $\vec\sigma = B^r$ for any
base block $B \in \Sigma(a, b)$. Then
\[
\Phi_{a, b, r} \mid C(y, \vec\sigma)
\]
as integers, regardless of the choice of $B$. Equivalently, every prime
power $p^e$ dividing $\Phi_{a, b, r}$ also divides $C(y, \vec\sigma)$.
\end{corollary}

\begin{proof}
Combining \eqref{eq:block-identity} and \eqref{eq:diff-factorisation},
\[
C(y, \vec\sigma) = C(b, B) \cdot \Phi_{a, b, r},
\]
which exhibits $\Phi_{a, b, r}$ as an integer factor of $C(y, \vec\sigma)$.
\end{proof}

\begin{corollary}[Recursive reduction]\label{cor:reduction}
Under the hypotheses of Lemma~\ref{lem:block}, suppose
$\gcd(2^a - 3^b, \Phi_{a, b, r}) = 1$. Then for $\vec\sigma = B^r$,
\[
(2^x - 3^y) \mid C(y, \vec\sigma)
\quad \Longleftrightarrow \quad
(2^a - 3^b) \mid C(b, B).
\]
\end{corollary}

\begin{proof}
By Corollary~\ref{cor:forced} and \eqref{eq:diff-factorisation},
$(2^x - 3^y) \mid C(y, \vec\sigma)$ if and only if
$(2^a - 3^b) \mid C(b, B) \cdot \Phi_{a, b, r}$. Coprimality forces
$(2^a - 3^b) \mid C(b, B)$.
\end{proof}

\begin{remark}[On the coprimality hypothesis]
Reducing $\Phi_{a, b, r}$ modulo $(2^a - 3^b)$, every $2^a$ is replaced by $3^b$,
so $\Phi_{a, b, r} \equiv \sum_{j=0}^{r-1} 3^{j b} \cdot 3^{(r-1-j) b}
= r \cdot 3^{(r-1) b} \pmod{2^a - 3^b}$. Hence
$\gcd(2^a - 3^b, \Phi_{a, b, r}) = \gcd(2^a - 3^b, r \cdot 3^{(r-1) b})$.
This is $1$ whenever $\gcd(2^a - 3^b, 6 r) = 1$, which holds in every case
empirically tested in this note.
\end{remark}


\section{The structure of the cyclic obstruction}

\begin{definition}[Periodic, sporadic, irreducible]\label{def:regimes}
Let $(x, y)$ be a parameter pair with $2^x > 3^y$ and $\gcd(x, y) = g$.
\begin{itemize}
\item A pattern $\vec\sigma \in \Sigma(x, y)$ is \emph{periodic of order $r$}
if $r \mid g$, $r \ge 2$, and $\vec\sigma = B^r$ for some base block
$B \in \Sigma(x/r, y/r)$.
\item A pattern with primitive period exactly $y$ is \emph{sporadic}.
\item The parameter pair $(x, y)$ is \emph{irreducible} if $g = 1$, and
\emph{reducible} if $g \ge 2$.
\end{itemize}
\end{definition}

\begin{conjecture}[Restricted zero-avoidance]\label{conj:main}
Let $(x, y)$ satisfy $2^x > 3^y$. If $\vec\sigma \in \Sigma(x, y)$ satisfies
$C(y, \vec\sigma) \equiv 0 \pmod{2^x - 3^y}$, then $(x, y) = (2 y, y)$
and $\vec\sigma$ is the trivial cycle.
\end{conjecture}

\begin{theorem}[Recursive structure of Conjecture~\ref{conj:main}]\label{thm:reduction-structure}
Suppose Conjecture~\ref{conj:main} holds for every irreducible parameter pair
with $2^a > 3^b$, and additionally that no sporadic pattern violates it.
Assume the coprimality $\gcd(2^a - 3^b, \Phi_{a, b, r}) = 1$ in all reductions.
Then Conjecture~\ref{conj:main} holds for every reducible parameter pair as well.
\end{theorem}

\begin{proof}
Let $(x, y)$ be reducible with $\gcd(x, y) = g \ge 2$, and let
$\vec\sigma \in \Sigma(x, y)$ satisfy $(2^x - 3^y) \mid C(y, \vec\sigma)$.
If $\vec\sigma$ is sporadic, the conclusion holds by hypothesis. Otherwise
$\vec\sigma = B^r$ for some $r \mid g$, $r \ge 2$, with $B \in \Sigma(x/r, y/r)$.
By Corollary~\ref{cor:reduction}, $(2^{x/r} - 3^{y/r}) \mid C(y/r, B)$. Iterating,
either $\vec\sigma$ remains periodic indefinitely (forcing $a_i = 2$ for all $i$
in the limit, hence the trivial cycle), or the recursion terminates at a sporadic
pattern of some reducible pair, or at an irreducible base case. By hypothesis,
neither of the latter two yields a non-trivial cycle.
\end{proof}

The remaining content of Conjecture~\ref{conj:main} is therefore concentrated in
two regimes: \emph{irreducible parameter pairs}, and \emph{sporadic patterns} of
reducible pairs. For both regimes, the present note offers only empirical evidence.


\section{Empirical observations}

\subsection{Aggregate avoidance}\label{sec:aggregate}

Across 22 parameter pairs with $2^x > 3^y$, exhaustive enumeration totalled
$1{,}769{,}706{,}522$ admissible patterns. The naive Crandall benchmark
$E_{\mathrm{Cr}}(x, y) = |\Sigma(x, y)| / (2^x - 3^y)$ aggregates to
$\sum E_{\mathrm{Cr}} \approx 7.21$ across non-trivial pairs. The observed
count of non-trivial integrality events is $0$.

\begin{table}[h]
\centering
\small
\caption{Selected enumerated pairs. The class column distinguishes
trivial ($x = 2 y$), reducible ($\gcd(x, y) = r \ge 2$, marked
\emph{red.\,r}), and irreducible (\emph{irr.}). Full 22-row table in the archive.}
\label{tab:agg}
\begin{tabular}{rrrrrrl}
\toprule
$y$ & $x$ & $|\Sigma|$ & $2^x - 3^y$ & $E_\mathrm{Cr}$ & non-triv.\ zeros & class \\
\midrule
5  & 8  & 35           & 13            & 2.69 & 0 & irr.\ \\
9  & 15 & 3{,}003      & 13{,}085      & 0.23 & 0 & red.\ 3 \\
10 & 16 & 5{,}005      & 6{,}487       & 0.77 & 0 & red.\ 2 \\
12 & 20 & 75{,}582     & 517{,}135     & 0.15 & 0 & red.\ 4 \\
13 & 22 & 293{,}930    & 2{,}599{,}981 & 0.11 & 0 & irr.\ \\
15 & 24 & 817{,}190    & 2{,}428{,}309 & 0.34 & 0 & red.\ 3 \\
16 & 26 & 3{,}268{,}760 & 24{,}062{,}143 & 0.14 & 0 & red.\ 2 \\
17 & 27 & 5{,}311{,}735 & 5{,}077{,}565 & 1.05 & 0 & irr.\ \\
19 & 31 & 86{,}493{,}225 & 985{,}222{,}181 & 0.09 & 0 & irr.\ \\
21 & 34 & 573{,}166{,}440 & 6{,}719{,}515{,}981 & 0.09 & 0 & irr.\ \\
22 & 35 & 927{,}983{,}760 & 2{,}978{,}678{,}759 & 0.31 & 0 & irr.\ \\
\midrule
\multicolumn{4}{r}{\textit{aggregate over all 22 enumerated non-trivial rows}}
& 7.21 & 0 & \\
\bottomrule
\end{tabular}
\end{table}

\subsection{Concrete instances of the block reduction}\label{sec:instances}

We illustrate Lemma~\ref{lem:block} on three reducible cases.

\paragraph{Case $(x, y) = (24, 15)$, $r = 3$, $(a, b) = (8, 5)$.}
Here
$2^{24} - 3^{15} = 13 \cdot 186{,}793$,
where $13 = 2^8 - 3^5$ and
$186{,}793 = 2^{16} + 2^8 \cdot 3^5 + 3^{10} = \Phi_{8, 5, 3}$.
For every base block $B \in \Sigma(8, 5)$ (there are 35 such blocks), the
identity
\[
C(15, B^3) = C(5, B) \cdot 186{,}793
\]
holds in $\mathbb{Z}$. Hence $186{,}793 \mid C(15, B^3)$ automatically, and
$13 \mid C(15, B^3) \iff 13 \mid C(5, B)$. Direct enumeration shows that
$C(5, B) \not\equiv 0 \pmod{13}$ for any of the 35 blocks. The 35 patterns
of primitive period $5$ in $\Sigma(24, 15)$ are therefore not solutions of
the cycle equation, by the lemma.

\paragraph{Case $(x, y) = (16, 10)$, $r = 2$, $(a, b) = (8, 5)$.}
Here
$2^{16} - 3^{10} = 6{,}487 = 13 \cdot 499$,
where $13 = 2^8 - 3^5$ and $499 = 2^8 + 3^5 = \Phi_{8, 5, 2}$. For every
$B \in \Sigma(8, 5)$,
\[
C(10, B^2) = C(5, B) \cdot 499.
\]
Hence $499 \mid C(10, B^2)$ automatically, and again
$13 \mid C(10, B^2) \iff 13 \mid C(5, B)$. None of the 35 blocks satisfies the
latter, so the periodic patterns of order $2$ in $\Sigma(16, 10)$ contribute
no cycle solutions.

\paragraph{Case $(x, y) = (26, 16)$, $r = 2$, $(a, b) = (13, 8)$.}
Here
$2^{26} - 3^{16} = (2^{13} - 3^8)(2^{13} + 3^8) = 1{,}631 \cdot 14{,}753$,
with $1{,}631 = 7 \cdot 233$ and $14{,}753$ prime. The block reduction with
$r = 2$ gives, for every $B \in \Sigma(13, 8)$ (there are $\binom{12}{7} = 792$
such blocks),
\[
C(16, B^2) = C(8, B) \cdot 14{,}753.
\]
Hence $14{,}753 \mid C(16, B^2)$ automatically, and
$1{,}631 \mid C(16, B^2) \iff 1{,}631 \mid C(8, B)$. Direct enumeration shows
that none of the 792 blocks satisfies $1{,}631 \mid C(8, B)$.

\subsection{The mutual exclusion phenomenon, revisited}\label{sec:mutex}

The mutual exclusion observation recorded in version 1 of this note now
admits a partial algebraic interpretation. For each reducible pair $(x, y)$
with $\gcd(x, y) = r \ge 2$, the prime divisors of $2^x - 3^y$ split into:
\begin{itemize}
\item primes dividing $\Phi_{a, b, r}$, which divide $C(y, B^r)$
\emph{automatically} (Corollary~\ref{cor:forced});
\item primes dividing $2^a - 3^b$, which divide $C(y, B^r)$ only via
$C(b, B)$ (Corollary~\ref{cor:reduction}).
\end{itemize}
On the periodic stratum, these two groups of primes act on disjoint
arithmetic conditions: the first via algebraic identity, the second via
the cycle equation at the smaller pair $(a, b)$. The empirical mutual
exclusion of divisibility events at distinct prime divisors of $2^x - 3^y$
is a direct consequence of this algebraic split, on the periodic stratum.

\subsection{Sporadic patterns}\label{sec:sporadic}

For each reducible pair tested, a small number of admissible patterns of
primitive period $y$ also satisfy divisibility by some prime factor of
$2^x - 3^y$, without admitting a block decomposition. We refer to these as
\emph{sporadic patterns}. For example, in $\Sigma(24, 15)$ the set
$A_{186{,}793} = \{\vec\sigma : 186{,}793 \mid C(15, \vec\sigma)\}$ contains
50 patterns: 35 are periodic of order $3$ (covered by
Lemma~\ref{lem:block}), and 15 form a single rotation class of primitive
period $15$, which is sporadic. None of the sporadic patterns observed in
any tested case extends to a full integer cycle.

The block reduction lemma does not apply to sporadic patterns. Their
arithmetic structure remains an open question.


\section{Limitations and open questions}

\begin{enumerate}
\item Lemma~\ref{lem:block} addresses only periodic patterns. The
conjecture restricted to sporadic patterns is unproved.

\item For irreducible parameter pairs $(\gcd(x, y) = 1)$ no block
reduction is available. The empirical evidence at irreducible pairs is
recorded in Table~\ref{tab:agg} but unsupported by the present mechanism.
The case $(8, 5)$ is the smallest irreducible base case ($|\Sigma| = 35$,
$2^8 - 3^5 = 13$); cycle integrality fails by direct enumeration but
without an algebraic explanation.

\item The recursive structure of
Theorem~\ref{thm:reduction-structure} terminates either at irreducible
pairs or at sporadic patterns of reducible pairs. Both are unresolved.
The conjecture for the cyclic component of Collatz is therefore
reformulated as the conjunction of two independent conjectures: one on
sporadic patterns, one on irreducible base cases.

\item The coprimality hypothesis
$\gcd(2^a - 3^b, \Phi_{a, b, r}) = 1$ holds in every case empirically
tested, and is implied by $\gcd(2^a - 3^b, 6 r) = 1$. We do not know
whether the hypothesis can fail in arithmetically singular cases.

\item Enumeration is bounded by $|\Sigma| \le 10^9$ on a single machine.
Convergent pairs of $\log_2 3$ beyond $(65, 41)$ are inaccessible.

\item The literature survey is preliminary. The block-reduction
identity \eqref{eq:block-identity} is elementary and may have appeared
in the cycle-equation literature; any precedent will be acknowledged.
The use of this identity to recursively reduce the cyclic obstruction,
to the author's knowledge, has not been recorded explicitly.

\item The note does not prove the absence of non-trivial Collatz cycles.
It identifies a partial mechanism, isolates the unsolved residue, and
relates it to known classical questions on $S$-unit sums.

\end{enumerate}


\paragraph{Code and data.}
A reproducibility archive containing the enumeration code, all parameter
data, the block-reduction verification scripts, and the per-case
prime-power decompositions accompanies the public version of this note.

\paragraph{Acknowledgements.}
This version corrects and substantially extends version 1
(Zenodo, May 2026) by isolating the algebraic mechanism behind the
empirical mutual exclusion observed there, and by introducing the
periodic / sporadic / irreducible classification of patterns.


\begin{thebibliography}{12}
\small

\bibitem{belaga1998}
E.~B.~Belaga and M.~Mignotte.
\emph{Embedding the $3x+1$ conjecture in a $3x+d$ context.}
Experimental Mathematics \textbf{7} (1998), 145--151.

\bibitem{belaga2003}
E.~B.~Belaga.
\emph{Effective polynomial upper bounds to perigees and numbers of
$(3x+d)$-cycles of a given oddlength.}
Acta Arithmetica \textbf{106} (2003), 197--206.

\bibitem{belaga2006}
E.~B.~Belaga and M.~Mignotte.
\emph{The Collatz Problem and Its Generalizations: Experimental Data.
Table 1. Primitive Cycles of $(3n+d)$-mappings.}
HAL hal-00129727, 2006.

\bibitem{bohm1978}
C.~B\"ohm and G.~Sontacchi.
\emph{On the existence of cycles of given length in integer sequences
like $x_{n+1} = x_n / 2$ if $x_n$ is even, and $x_{n+1} = 3 x_n + 1$
otherwise.}
Atti della Accademia Nazionale dei Lincei, Rendiconti \textbf{64}
(1978), 260--264.

\bibitem{crandall1978}
R.~E.~Crandall.
\emph{On the ``$3x+1$'' problem.}
Mathematics of Computation \textbf{32} (1978), 1281--1292.

\bibitem{dhiman2026}
M.~Dhiman and R.~Pandey.
\emph{2-Adic Obstructions to Presburger-Definable Characterizations of
Collatz Cycles.}
arXiv:2601.12772, 2026.

\bibitem{eliahou1993}
S.~Eliahou.
\emph{The $3x+1$ problem: new lower bounds on nontrivial cycle lengths.}
Discrete Mathematics \textbf{118} (1993), 45--56.

\bibitem{hercher2022}
C.~Hercher.
\emph{There are no Collatz $m$-cycles with $m \le 91$.}
arXiv:2201.00406, 2022.

\bibitem{holden2011}
D.~Holden.
\emph{Results on the $3x+1$ and $3x+d$ conjectures.}
Fibonacci Quarterly \textbf{49}(2) (2011), 131--133.

\bibitem{lagarias1985}
J.~C.~Lagarias.
\emph{The $3x+1$ problem and its generalizations.}
American Mathematical Monthly \textbf{92} (1985), 3--23.

\bibitem{sdw2005}
J.~Simons and B.~de Weger.
\emph{Theoretical and computational bounds for $m$-cycles of the $3n+1$
problem.}
Acta Arithmetica \textbf{117} (2005), 51--70.

\bibitem{simons2008}
J.~L.~Simons.
\emph{On the (non-)existence of $m$-cycles for generalized Syracuse
sequences.}
Acta Arithmetica \textbf{131} (2008), 217--254.

\bibitem{steiner1977}
R.~P.~Steiner.
\emph{A theorem on the Syracuse problem.}
Proceedings of the 7th Manitoba Conference on Numerical Mathematics
(1977), 553--559.

\end{thebibliography}

\end{document}
